\section{Introduction}
\label{sec:introduction}

\subsection{Motivation}

When a malware sample reaches an incident response or threat
intelligence team, the useful output is rarely the binary itself: it
is a structured, actionable summary of which indicators of compromise
(IOCs) it exhibits, which MITRE ATT\&CK~\cite{mitre-attack} techniques
it implements, and artifacts ready to feed detection engineering
(Sigma rules, YARA signatures) or a threat intelligence platform such
as MISP~\cite{misp}. A quick first pass at that summary is often
already available: a hash lookup against a reputation aggregator such
as VirusTotal returns vendor verdicts and sandbox behavioral tags
within seconds. What that first pass does not provide is a
\emph{reconciled} answer, since the many signals it surfaces are
displayed side by side rather than checked against each other, so
turning them into a defensible technique mapping for one specific
sample usually still means a human analyst manually detonating the
sample in a sandbox, tracing its code in a disassembler, and writing a
report by hand. That process is thorough but slow, and it does not
scale to the volume of samples a security operations team realistically
encounters.

Automated sandboxes such as CAPE~\cite{cape} already produce large
volumes of raw behavioral data: API call traces, dropped files, network
activity, and community-maintained signature matches. Static analysis
tools extract imports, strings, and packed-payload artifacts directly
from the binary. What is generally missing is the layer that sits
between these raw findings and something an analyst or a downstream
platform can act on directly: reconciling static and dynamic evidence
for the \emph{same} claim, mapping that evidence onto a standard
technique taxonomy with a defensible confidence level, and exporting it
in a format threat intelligence platforms already consume.

\subsection{Problem Statement}

Three specific problems motivate this work, each identified while
building and validating the pipeline described in this thesis:

\begin{enumerate}
  \item \textbf{Confidence reconciliation across sources.} A technique
    observed by both static analysis (an import combination) and
    dynamic analysis (a sandbox signature) is stronger evidence than
    either alone, but only if the two observations genuinely refer to
    the same technique. A naive implementation that compares
    observations by exact technique identifier misses the common case
    where one source reports a parent technique and the other reports
    a specific sub-technique for what is, in reality, the same
    behavior.
  \item \textbf{Multi-stage malware.} A loader that drops a rootkit
    driver, an AV-killer module, and a remote-access client has real
    capabilities distributed across several binaries, not concentrated
    in the one file initially detonated. Naively merging every dropped
    component's evidence into the parent sample's own confidence score
    would misattribute capability, since the dropped payload's own
    behavior is not proof the loader itself performs that behavior.
  \item \textbf{Trusting the output.} A pipeline that reports high
    confidence for a wrong finding is worse than one that reports
    nothing, because it actively misleads a downstream consumer. This
    motivates treating every detection rule as a claim that needs
    evidence, and every evaluation result as something to audit before
    trusting, not just a number to report.
\end{enumerate}

\subsection{Contributions}

This thesis makes the following contributions, each substantiated in
detail in later sections:

\begin{itemize}
  \item A \textbf{cross-source and cross-level confidence
    reconciliation model} (\S\ref{sec:reconciliation}) that extends
    the standard ``both sources agree $\Rightarrow$ high confidence''
    rule to also recognize agreement between a parent technique
    reported by one source and its specific sub-technique reported by
    the other, a gap confirmed present in real evaluation data, not
    merely theoretical.
  \item A \textbf{resubmission loop} (\S\ref{sec:resubmission}) that
    independently analyzes every file a multi-stage sample drops or
    injects during detonation, tagging each with lineage back to the
    parent sample while keeping its evidence out of the parent's own
    confidence-scored output.
  \item A \textbf{reproducible, evidence-sourced static detection rule
    set} (\S\ref{sec:coverage}), extended systematically from roughly
    30 to approximately 85 Windows-relevant MITRE ATT\&CK technique
    identifiers, with every addition traced to the technique's own
    official MITRE description and verified against the validated
    sample set before being trusted as non-overfit.
  \item An \textbf{evaluation methodology and audit process}
    (\S\ref{sec:evaluation}) that goes beyond precision/recall:
    confidence-tier- and source-agreement-stratified precision to test
    whether the reconciliation model's own claims hold, and a
    systematic false-positive audit in which every discrepancy between
    the pipeline's output and manually-verified ground truth was
    individually diagnosed to a specific, verifiable cause before
    being fixed, left open, or attributed to a ground-truth gap.
\end{itemize}

\subsection{Document Structure}

Section~\ref{sec:related-work} situates this work relative to existing
sandboxing tools, static ATT\&CK-mapping approaches, and reputation
aggregators such as VirusTotal. Section~\ref{sec:architecture}
describes the pipeline's architecture end to end. Section~\ref{sec:evaluation}
presents the evaluation methodology and results, including the
false-positive audit. Section~\ref{sec:case-studies} details four
concrete case studies, including a real architectural bug found and
fixed through the evaluation process, that illustrate the project's
validation discipline more concretely than aggregate numbers can.
Section~\ref{sec:limitations} states the pipeline's limitations
precisely, each tied to a specific diagnosed cause.
Section~\ref{sec:future-work} outlines concrete extensions, including
one, Linux/ELF dynamic analysis for IoT botnet families, evaluated and
deliberately deferred for reasons given explicitly rather than left
implicit. Section~\ref{sec:conclusion} concludes.
Appendix~\ref{app:rules} tabulates the detection rule set added during
coverage extension. The full pipeline source code is published in the
project's public repository rather than reproduced here.
